% ============================================================
\section{Conclusiones}
% ============================================================

La evidencia recogida en los tres benchmarks permite concluir lo siguiente:

\begin{enumerate}
  \item La GPU aceleró todas las cargas evaluadas, pero con magnitudes muy
        distintas según la naturaleza del cómputo: 2.4--2.7$\times$ en el
        entrenamiento de YOLO26n-seg (limitado por la T4 y el pipeline de
        datos de 2 vCPU), 3.2--4.4$\times$ en inferencia profunda sobre video
        (YOLO y Real-ESRGAN en la RTX 3070) y 2.9$\times$ en el pipeline
        clásico de bordes en C++. El factor determinante es la razón entre
        cómputo paralelizable y latencia de transferencia CPU$\leftrightarrow$VRAM.
  \item El fine-tuning convergió correctamente: la pérdida de clasificación
        cayó de 4.503 a 0.401 entre las épocas 1 y 60 y el mAP50 de cajas
        subió de 0.192 a 0.792, con un consumo estable de VRAM de
        3.31--4.6\,GiB, muy por debajo de los 15\,GiB de la T4.
  \item En inferencia GPU ambas variantes YOLO convergieron a
        $\approx$122--125 FPS pese a rendir distinto en CPU (29.3 vs.\ 38.9
        FPS), lo que demuestra que, una vez acelerada la red, el cuello de
        botella se traslada a las etapas secuenciales del bucle (decodificación
        y visualización), en concordancia con la ley de Amdahl.
  \item El costo sistémico de forzar cargas profundas en CPU es severo:
        Real-ESRGAN llevó al i9-12900K a 91\,$^{\circ}$C con 52\,\% de uso
        agregado para entregar 0.36 FPS, mientras la GPU resolvió la misma
        carga a 1.57 FPS con la CPU al 5\,\%. La GPU no solo es más rápida:
        libera al resto del sistema.
  \item En procesamiento clásico, la CPU ya supera el tiempo real (150 FPS a
        480p), por lo que la aceleración CUDA (432 FPS, 2\,\% de utilización,
        172\,MiB de VRAM) solo se justifica en resoluciones mayores, múltiples
        flujos o pipelines más profundos; el diseño de una única subida y una
        única bajada por cuadro fue la condición necesaria para que la GPU
        ganara en este caso.
\end{enumerate}
